% ---------------------------------------------------------------------
%  Chapter 12 : Phase J --- The Temporal Loop-Integral Engine
% ---------------------------------------------------------------------
\chapter{Phase J: The Temporal Loop-Integral Engine}\label{ch:phasej-temporal}

Every chapter so far has been preparation for this one. The theory file gave us
an action; the field-theory core expanded it into vertices; the propagator
chapter built the free two-point function; the mean-field solve found the
saddle; the enumeration and type-assignment chapters produced a list of typed
Feynman diagrams, each tagged with a scalar prefactor. What we do \emph{not} yet
have is a single number. A diagram is still an abstract graph plus a
prefactor; turning it into a value --- a function of the external times we can
plot against a simulation --- is the job of \term{Phase~J}, the final,
time-domain evaluation layer of the temporal pipeline.

This is the numerical heart of the whole temporal branch. It is also the most
intricate piece of code in the repository: \file{final\_integral.py} is a single
file of roughly 5{,}300 lines. The intricacy is not gratuitous. The integral a
diagram represents is, in general, a multi-dimensional integral over a region
carved out by causality constraints, and the engine goes to considerable
lengths to evaluate that integral \emph{analytically} --- in closed form ---
rather than handing it to a black-box numerical quadrature routine. This chapter
explains, from the ground up, why that integral has a closed form, what the
closed form is in each of the cases the engine distinguishes, and how the code
realises it.

\begin{note}
The name ``Phase~J'' is a historical artifact. The temporal pipeline was
originally built as a hybrid: the early phases (``I'' and before) work in
frequency space, and ``J'' was the last phase, the one that returns to the time
domain to finish the job. It began life as a tree-level ``minimum viable
product,'' but the core integrator (\code{integrate\_diagram}) has since been
generalised to handle \emph{any} loop order. The label stuck even though the
algorithm outgrew its original scope.
\end{note}

\section{What Phase J computes}

In the Martin--Siggia--Rose--Janssen--De~Dominicis (\msrjd) formalism the
quantity we ultimately want is a time-domain correlation or response function of
a stochastic dynamical system --- for example the autocorrelation
$C(\tau) = \avg{\delta n(t)\,\delta n(t+\tau)}$ of a Hawkes process, or a
higher-order cumulant of any of the temporal models in the library. Perturbation
theory expresses that correlator as a sum over Feynman diagrams. Each diagram is
a graph whose

\begin{itemize}
  \item \textbf{edges} are \emph{retarded propagators} $\Gret(t_v - t_u)$ --- the
        linear response of one field to another, carrying a Heaviside step
        $\Theta$ that enforces causality (a response can only happen \emph{after}
        its cause), and whose
  \item \textbf{internal vertices} are interaction points whose \emph{times}
        must be integrated over the whole real line, subject to the causal
        orderings the Heavisides impose.
\end{itemize}

So the value of one diagram is a multi-dimensional integral of the schematic
form
\begin{equation}
  C_\Gamma(\text{external times})
   \;=\; \text{prefactor} \cdot
   \int \dd s_1 \cdots \dd s_m \;\prod_{e \in E} \Gret\big(t_{\text{head}(e)} - t_{\text{tail}(e)}\big),
  \label{eq:diagram-integral}
\end{equation}
where $s_1, \dots, s_m$ are the internal-vertex times. \textbf{Phase~J is the
machine that performs that integral}, for every diagram, at every loop order,
returning a Python callable \code{f(*external\_times) -> complex} that you can
evaluate on any grid of external times.

\begin{defn}
A \textbf{retarded propagator} $\Gret(\Delta t)$ is the linear response of a
physical field to a perturbation in its conjugate response field a time
$\Delta t$ earlier. ``Retarded'' means it vanishes for $\Delta t < 0$: the
effect cannot precede the cause. In the diagrams this causality is carried by an
explicit Heaviside factor $\Theta(\Delta t)$ on every edge, and those Heavisides
are precisely what turn \eqref{eq:diagram-integral} from an integral over all of
$\mathbb{R}^m$ into an integral over a bounded \emph{causal region}.
\end{defn}

The cleverness of the engine is that it almost never falls back to brute-force
numerical quadrature. Because every retarded propagator turns out to be a
\emph{finite sum of decaying complex exponentials} (we derive this in
\S\ref{sec:modesum}), the whole integrand is a sum of terms of the form
$A \cdot \exp(\text{linear combination of the times})$, and the integral of such
a term over the causal region --- which is a \emph{polytope}, a region bounded by
flat faces --- has a closed form. The engine therefore carries three dedicated
\emph{analytic} integrators, one each for $1$, $2$, and $\ge 3$ internal-vertex
integration variables, and only routes to \code{scipy}'s numerical quadrature
when the analytic path hits a genuine degeneracy or a floating-point overflow.

\subsection{Where it sits, and what feeds it}

Phase~J consumes four things, all produced by earlier phases:

\begin{enumerate}
  \item a list of \term{typed diagrams} (\code{TypedDiagram} objects) from the
        enumeration, typing, and symmetry-deduplication pipeline;
  \item a \emph{scalar prefactor} per diagram --- the combinatorial and coupling
        weight, from \code{classify\_coefficient\_factors};
  \item the retarded propagator in \term{pole--residue form}: a dictionary
        \code{propagator\_data} with keys \code{pole\_vals} (the poles $p_\alpha$),
        \code{C\_mats} (one residue matrix per pole), and \code{D\_delta} (the
        instantaneous $\delta$-function part);
  \item the user's external-field list and the numerical parameter
        substitutions \code{num\_params}.
\end{enumerate}

\noindent The public entry point is \code{compute\_correction\_td} in
\file{pipeline.py}, which loops over the diagrams calling \code{integrate\_diagram}
(in \file{final\_integral.py}) on each. Downstream, the result feeds the
notebook-facing callables \code{total\_C} and \code{total\_C\_batch} that produce
the theory curves you overlay on simulation data.

The three source files this chapter dissects are:

\begin{description}
  \item[\file{final\_integral.py}] the $\sim$5{,}300-line core: the per-diagram
        integrator and all of the analytic primitives.
  \item[\file{pipeline.py}] the orchestration and batching layer: it loops over
        diagrams and fans the $\tau$-grid evaluation out over worker processes.
  \item[\file{subgraph.py}] a near-stub for a frequency-domain loop-reduction
        route that was designed but never finished --- the time-domain integrator
        superseded it. We cover it briefly at the end so you know it is dead
        weight, not a live path.
\end{description}

\section{The retarded propagator as a mode sum}\label{sec:modesum}

Everything downstream rests on one fact: the time-domain retarded propagator is
a finite sum of complex exponentials. Let us derive it, because the code is a
literal transcription of the result.

\subsection{From frequency to time by residues}

For a \emph{linear} response system the frequency-domain propagator
$\widehat{G}(\omega)$ is a rational matrix function of $\omega$ --- a ratio of
polynomials, entry by entry. Its inverse Fourier transform, under the project's
fixed convention
\begin{equation}
  G(t) \;=\; \frac{1}{2\pi}\int \dd\omega \; \ee^{\ii \omega t}\,\widehat{G}(\omega),
  \label{eq:fourier-convention}
\end{equation}
is computed by closing the $\omega$ contour and summing residues. \emph{Why
upward?} For $t>0$ the kernel $\ee^{\ii\omega t}$ decays as
$\operatorname{Im}\omega\to+\infty$, so by Jordan's lemma the contour may be
closed by a large semicircle in the \emph{upper} half plane at no added cost,
enclosing exactly the upper-half-plane poles (and for $t<0$ one closes downward,
enclosing none). The causality
filter built into the propagator construction guarantees that every pole
$p_\alpha$ of $\widehat{G}$ lies in the \emph{upper} half plane,
$\operatorname{Im} p_\alpha > 0$. Each residue term therefore carries a factor
$\ee^{\ii p_\alpha t}$, which \emph{decays} for $t > 0$ (its exponent has real
part $-\operatorname{Im} p_\alpha < 0$) and \emph{grows} for $t < 0$. The
physical retarded propagator is this analytic sum multiplied by a Heaviside that
kills the growing $t<0$ half:
\begin{equation}
  \Gret[p, r](\Delta t)
   \;=\; \underbrace{\texttt{delta\_coeff}[p,r]\cdot\delta(\Delta t)}_{\text{instantaneous part}}
   \;+\; \Theta(\Delta t)\cdot\sum_\alpha C_\alpha[p,r]\,\ee^{\ii p_\alpha \Delta t}.
  \label{eq:GR-decomp}
\end{equation}

Three pieces of nomenclature, each of which has a direct counterpart in the
code:

\begin{itemize}
  \item The indices $[p, r]$ are the \textbf{(physical-row, response-column)}
        matrix indices. The convention is $G[p_i, r_i] = \avg{\phi_{p_i}\,
        \tilde n_{r_i}}$ --- the response of physical field $p_i$ to a
        response-field (hatted) source $r_i$. There is a transpose to watch
        here: the kernel matrix $K$ from which the propagator is built is laid
        out \code{[resp, phys]}, so to read off a retarded entry you index
        $G$\code{[phys, resp]}. The propagator chapter covers this; Phase~J just
        respects it.
  \item $\texttt{delta\_coeff}[p,r] = \lim_{\omega\to\infty}\widehat{G}[p,r](\omega)$
        is the \textbf{instantaneous response} --- the polynomial (non-proper)
        part of the rational $\widehat{G}$. It is nonzero exactly when there is
        an instantaneous coupling: for example the $\tilde n \times \delta n$
        term in an \msrjd{} action, where a $\tilde n$ source at time $t$
        produces an \emph{immediate} $\delta n$ at the same $t$. It is captured
        in \code{propagator\_data['D\_delta']}.
  \item $C_\alpha[p,r]$ is the residue of $\widehat{G}$ at pole $p_\alpha$,
        position $[p,r]$.
\end{itemize}

To keep the inner loops in plain Python arithmetic, the code stores not the pole
$p_\alpha$ but the combination $\lambda_\alpha := \ii\, p_\alpha$, so that a mode
reads simply $C_\alpha\,\ee^{\lambda_\alpha \Delta t}$ with no factor of $\ii$ to
re-apply on every evaluation.

\subsection{The \code{EdgeModeSum} record}

This entire object --- one edge's worth of \eqref{eq:GR-decomp} --- is packaged
into a single frozen dataclass, \code{EdgeModeSum}
(\file{final\_integral.py:116}). Its documented fields are the residue and pole
indices, the $\delta$-coefficient, and the list of modes:

\begin{lstlisting}
@dataclass(frozen=True)
class EdgeModeSum:
    ri: int
    pi: int
    delta_coeff: complex
    modes: tuple                 # tuple[tuple[complex, complex], ...]  i.e. (C_alpha, lambda_alpha)
    dt_c0: float = 0.0
    dt_int_pairs: tuple = ()     # sparse coefficients on integration vars
    dt_ext_pairs: tuple = ()     # sparse coefficients on free external times
\end{lstlisting}

\begin{note}
\code{@dataclass} is a standard-library decorator (from the \code{dataclasses}
module) that auto-generates a class's constructor, \code{\_\_repr\_\_}, and
equality from a list of typed field declarations --- you write the fields, Python
writes the boilerplate. \code{frozen=True} makes instances immutable (assigning
to a field raises), which is what lets these records be safely cached and shared
across the evaluation without fear of one code path mutating another's data.
\end{note}

The \code{modes} tuple holds the $(C_\alpha, \lambda_\alpha)$ pairs. The three
\code{dt\_*} fields are left empty at construction; they hold the linear form for
$\Delta t$ in terms of the integration variables and external times, and they
are filled in later, once we know which subset of edges we are in (see
\S\ref{sec:polytope}). Carrying them on the same object means the one
\code{EdgeModeSum} threads through the whole evaluation chain.

The builder \code{\_build\_edge\_mode\_sums(edge\_info, propagator\_data)}
(\file{final\_integral.py:150}) does the extraction \emph{once} per edge. It
reads the poles and residue matrices out of \code{propagator\_data} and converts
them to machine-precision Python \code{complex} numbers:

\begin{lstlisting}
modes_lambdas = tuple(complex(CDF(SR(p))) * 1j for p in pole_vals)
# ... then per edge, per pole:
residues = tuple(
    complex(CDF(SR(C_mats[k][pi, ri])))
    for k in range(n_poles)
)
modes = tuple(zip(residues, modes_lambdas))
\end{lstlisting}

\noindent The idiom \code{complex(CDF(SR(expr)))} is how this codebase turns a
fully-substituted symbolic scalar into a Python \code{complex}: \code{SR} coerces
to a symbolic expression, \code{CDF} (the Complex Double Field) evaluates it to
machine precision, and \code{complex(...)} hands back a native Python number. We
explain \code{SR} and \code{CDF} in \S\ref{sec:tools}. If the propagator data is
incomplete --- missing \code{pole\_vals} or \code{C\_mats} --- the builder returns
\code{None}, which is the signal for the caller to fall back to the slow
symbolic path.

\begin{gotcha}
Extracting the residues once, here, rather than re-extracting them inside the
per-call evaluator, is a deliberate and important optimisation. The original
code re-derived the smooth-part data on every \code{(diagram, subset)} call even
though it depends only on the $(p_i, r_i)$ of the edge. Lifting the extraction to
this single per-edge step is what keeps the hot numerical loop in plain
\code{cmath}, never touching Sage --- Sage is far too slow to call per
evaluation.
\end{gotcha}

\section{From a diagram to a polytope integral}\label{sec:polytope}

We now assemble \eqref{eq:diagram-integral} concretely and watch it collapse to
the analytic form. There are three moving parts: the $\delta$-subset expansion,
the elimination of variables by $\delta$-edge equations, and the half-space
constraints that build the polytope.

\subsection{The $\delta$-subset expansion}

A diagram with edge set $E$ has integrand $\prod_{e\in E}\Gret(\Delta t_e)$.
Substitute the two-piece decomposition \eqref{eq:GR-decomp} into each factor and
\emph{expand the product}. Each edge is taken either in its $\delta$-form or in
its smooth (exponential) form, so the product becomes a sum over subsets
$S\subseteq E$ of ``which edges are $\delta$'':
\begin{equation}
  C_\Gamma = \sum_{S\subseteq E}\int \dd s_1\cdots \dd s_m\;
    \underbrace{\Big(\prod_{e\in S}\texttt{delta\_coeff}[e]\,\delta(\Delta t_e)\Big)}_{\delta\text{ edges}}
    \cdot
    \underbrace{\Big(\prod_{e\notin S}\Theta(\Delta t_e)\sum_\alpha C_\alpha[e]\,\ee^{\lambda_\alpha \Delta t_e}\Big)}_{\text{smooth edges}}
    \cdot \text{prefactor}.
  \label{eq:subset-sum}
\end{equation}

This is the central organising idea of the engine. Within each subset:

\begin{itemize}
  \item A \textbf{$\delta$-edge equation} $\Delta t_e = 0$ (for $e\in S$) is
        linear in the vertex times. Each one is solved to \emph{eliminate one
        integration variable} by substitution --- it pins one internal vertex's
        time to another's. This \emph{reduces} the dimension $m$ of the remaining
        integral.
  \item If a $\delta$-edge equation has \emph{no integration variable left to
        solve for}, it becomes a constraint among the external times alone --- a
        \textbf{shot-noise $\delta(\tau)$ spike}. The correlator then contains a
        term $\propto \delta(a\cdot\tau + c)$. These are distributional and are
        \emph{not} added to the smooth callable; they are emitted separately as
        structured \code{delta\_contributions} (covered in
        \S\ref{sec:deltaspike}).
  \item A \textbf{smooth edge} contributes its mode sum
        $\sum_\alpha C_\alpha\,\ee^{\lambda_\alpha \Delta t_e}$ and, via its
        Heaviside, a \textbf{half-space constraint} $\Delta t_e > 0$. Together,
        the smooth edges' constraints carve a convex \term{polytope} out of the
        space of remaining integration variables.
\end{itemize}

\begin{defn}
A \textbf{polytope} is the intersection of finitely many half-spaces --- a
convex region with flat faces (a polygon in 2D, a polyhedron in 3D, and so on).
In Phase~J the polytope is the set of internal-vertex times that respect all the
causal orderings $\{\Delta t_e > 0\}$ a subset imposes. The integral
\eqref{eq:subset-sum} is, per subset, an exponential integrated over this
polytope.
\end{defn}

\begin{gotcha}
Naively the subset sum has $2^{|E|}$ terms. The engine shrinks this
dramatically: edges whose $\delta$-coefficient is numerically zero
(\code{|delta\_coeff| < 1e-15}) are \emph{forced smooth} --- they have no
$\delta$-form to branch on --- so the sum runs only over the edges that actually
carry an instantaneous part, turning $2^{|E|}$ into $2^{|\text{branch}|}$. For
most temporal models only a handful of edges (the noise vertices) have a
$\delta$-part at all, so $|\text{branch}|$ is small.
\end{gotcha}

\subsection{Affine $\Delta t$ and the pole-tuple expansion}

After $\delta$-elimination, every surviving $\Delta t_e$ is an \emph{affine}
(linear-plus-constant) function of the integration variables
$s = (s_0, \dots, s_{m-1})$, the free external times $t_{\text{free}}$, and a
constant:
\begin{equation}
  \Delta t_e = c^{(e)}_0 + \sum_v a^{(e)}_{\text{int}}[v]\, s_v + \sum_j a^{(e)}_{\text{ext}}[j]\, t_{\text{free}}[j].
  \label{eq:dt-affine}
\end{equation}

Now choose one pole $\alpha_e$ per smooth edge --- a \term{pole tuple}. With that
choice fixed, the product of the chosen modes is a \emph{single} exponential
$C_{\text{prod}}\cdot\exp\!\big(\sum_e \lambda_{\alpha_e}\Delta t_e\big)$. Group
the exponent by what it multiplies:
\begin{align}
  \alpha_s[v]  &= \sum_e \lambda_{\alpha_e}\, a^{(e)}_{\text{int}}[v]   &&\text{(coefficient of } s_v\text{)}, \label{eq:alpha-s}\\
  \gamma       &= \sum_e \lambda_{\alpha_e}\Big(c^{(e)}_0 + \sum_j a^{(e)}_{\text{ext}}[j]\,t_{\text{free}}[j]\Big) &&\text{(the } s\text{-independent part)}. \label{eq:gamma}
\end{align}
So the integrand for one pole tuple is
$\text{prefactor}\cdot C_{\text{prod}}\cdot \ee^{\gamma}\cdot \exp\!\big(\sum_v \alpha_s[v]\,s_v\big)$,
and the whole subset integral is the sum, over all
$\prod_e (\#\text{poles}_e)$ pole tuples, of
\begin{equation}
  \int_{\text{polytope}} \exp\!\Big(\sum_v \alpha_s[v]\,s_v\Big)\,\dd s.
  \label{eq:exp-over-polytope}
\end{equation}
Each of these is a closed-form \emph{exponential-over-a-polytope}. This is the
integral the three analytic backends evaluate.

\begin{defn}
A \textbf{pole tuple} is a choice of exactly one pole $\alpha_e$ for each smooth
edge $e$ of a subset. Because the propagator's time-domain form is a sum over
poles, the product over edges of those sums expands into a sum over pole tuples;
each tuple contributes one pure exponential, and the polytope integral of a pure
exponential is closed-form. The number of pole tuples is the product of the
per-edge pole counts.
\end{defn}

\subsection{The plan cache}

The coefficients $\alpha_s[v]$ depend only on the pole tuple and the diagram
geometry --- \emph{not} on the value of the external times. Only $\gamma$ depends
on $t_{\text{free}}$, and it does so \emph{linearly}. The engine exploits this by
pre-computing, once per subset, everything that is $\tau$-invariant. That is what
\code{\_build\_modesum\_plan(...)} (\file{final\_integral.py:595}) does: it
returns a dict whose keys are \code{pole\_tuples}, \code{alphas\_per\_tuple} (the
$\alpha_s[v]$ of \eqref{eq:alpha-s}), \code{gamma\_const\_per\_tuple} ($\gamma$
evaluated at $t_{\text{free}}=0$), and \code{gamma\_slope\_per\_tuple\_per\_ext}
($\partial\gamma/\partial t_{\text{free}}[j]$). The per-$\tau$ work then collapses
to a single linear combination
\[
  \gamma = \gamma_{\text{const}} + \sum_j \text{slope}[j]\cdot t_{\text{free}}[j],
\]
followed by the closed-form polytope integral. The expensive part --- enumerating
pole tuples and building the $\alpha_s$ vectors --- is paid once, not once per
grid point.

The pole tuples themselves are produced by \code{\_enumerate\_pole\_tuples}
(\file{final\_integral.py:529}), a generator yielding $(C_{\text{product}},
\lambda\text{s})$ over the Cartesian product of the per-edge modes. It advances a
stack of indices by hand rather than calling \code{itertools.product}, to avoid
the per-tuple object-allocation overhead; an empty edge set yields the single
trivial tuple $(1.0+0j, ())$.

\section{Translation invariance and external-Wick sums}\label{sec:wick}

Before we reach the integrators, two correctness subtleties have to be handled in
the construction of the callable: pinning the origin, and summing over Wick
contractions of identical external legs.

\subsection{Origin pinning}

A stationary correlator depends only on \emph{time differences}. There is no
absolute clock, so we are free to fix one external leaf at $t=0$. The engine pins
the leaf at canonical position \code{origin\_leaf\_idx} (default $0$) to
\code{SR(0)} during integrand construction. The callable it returns therefore
computes a function of $(t_j - t_{\text{origin}})$ --- exactly the physical
content, with the redundant overall time shift removed. For a $k=2$
autocorrelator this means \code{total\_C(t\_1, t\_2)} returns $C$ at
$\tau = t_2 - t_1$.

\subsection{Summing Wick contractions, and the compensation divisor}

When several external legs carry the \emph{same} field type --- say two $\delta
n_1$ legs --- the symmetry-deduplication upstream has merged diagrams that differ
only in which same-type leaf attaches where. Phase~J must compensate by summing
over all such leaf permutations. Concretely, it enumerates every
field-respecting assignment of canonical positions to leaves (the
\code{\_all\_mappings} list), sums the integrand over them, and divides by the
\textbf{external-Wick compensation factor}
\begin{equation}
  \texttt{comp} \;=\; \frac{|\Aut(\Gamma,\ \text{leaves free})|}{|\Aut(\Gamma,\ \text{leaves fixed})|}.
  \label{eq:wick-comp}
\end{equation}
This is computed by \code{external\_wick\_compensation} in
\file{symmetry.py:343}. The justification is orbit--stabilizer: two assignments
yield the same pinned-external diagram iff they differ by an automorphism that
permutes leaves, and the stabilizer of any assignment is exactly
$\Aut_{\text{fixed}}$, so the mapping sum reproduces each distinct pinned diagram
$|\Aut_{\text{free}}|/|\Aut_{\text{fixed}}|$ times. Dividing by that index makes
the sum equal the sum over \emph{distinct} pinned diagrams --- which is what the
Feynman rules require, since each pinned diagram already carries its full weight
through its combinatorial factor (whose denominator is the same
$\Aut_{\text{fixed}}$).

\begin{gotcha}
The compensation divisor must be the \emph{exact} orbit--stabilizer index of
\eqref{eq:wick-comp}, not a cheaper proxy. A previous implementation approximated
it as $\prod N!$ over leaves grouped by their attachment vertex's
\emph{role signature} (a shallow depth-1 invariant). That over-divides whenever
two attachment vertices share a role signature but are \emph{not} related by any
automorphism --- e.g. the $k=4$ OU$+\varepsilon x^3$ one-loop cascades, where
three noise vertices all read ``noise feeding a cubic'' yet hang off
structurally distinct cubics. The resulting $\times\frac13$ / $\times\frac12$
deficits broke every $k\ge 3$ one-loop cumulant while coincidentally leaving
$k=2$ exact. The fix (validated against the exact Boltzmann series
$\kappa_4 = -6\varepsilon + 126\varepsilon^2$) is the true automorphism index.
When \code{external\_fields} is not supplied and a field-count mismatch forces
the identity-mapping fallback, the divisor must be exactly $1$, never the old
heuristic.
\end{gotcha}

\subsection{Re-pinning the origin per permutation}

There is a second, subtler point in the final callable
(\file{final\_integral.py:3759}). The integrand was \emph{built} with the origin
leaf pinned to $0$. When a non-identity permutation relabels which leaf occupies
which canonical position, the origin leaf may land at a non-zero time in the
user's input frame. The callable therefore time-translates each permuted
configuration so the origin returns to $0$ --- it subtracts the permuted origin's
time from all the others:

\begin{lstlisting}
total = 0.0 + 0.0j
for perm in _perms:
    permuted = [ext_time_values[perm[j]] for j in range(_k)]
    if origin_leaf_idx is not None:
        t_origin = permuted[origin_leaf_idx]
        # ... subtract t_origin from the free legs ...
\end{lstlisting}

\begin{gotcha}
This per-permutation re-pinning is load-bearing. For an \emph{asymmetric}
integrand --- which is every one-loop-and-higher diagram with a distinguished leaf,
such as a cubic-vertex tadpole --- the different permutations are physically
distinct evaluations (for a $k=2$ swap, $V(\tau)$ and $V(-\tau)$) and must be
genuinely summed, not merely counted. Before this fix the swap permutation was
fed the pinned origin's actual time, $\code{free\_val}=0$, producing spurious
asymmetry in $C(\tau)$ for every non-tree $k=2$ identical-externals case.
\end{gotcha}

\section{The $m=1$ interval integrator}\label{sec:m1}

We now turn to the three analytic backends, in order of increasing dimension.
The dispatch on $m$ (the number of integration variables surviving
$\delta$-elimination) happens inside each subset's contribution closure.

For $m=1$ the polytope is an interval $[L, U]$ on the single integration variable
$s$. With one variable, \eqref{eq:exp-over-polytope} is the elementary
\begin{equation}
  \int_L^U \ee^{\alpha_s\, s + \gamma}\,\dd s
   = \ee^{\gamma}\cdot\frac{\ee^{\alpha_s U} - \ee^{\alpha_s L}}{\alpha_s},
  \label{eq:m1-closed}
\end{equation}
with the removable-singularity limit $(U-L)\,\ee^{\gamma}$ as $\alpha_s\to 0$. The
implementation is \code{\_integrate\_1d\_polytope\_modesum}
(\file{final\_integral.py:1976}). The bounds $L$ and $U$ come from intersecting
the smooth-edge half-line constraints $a^{(e)}_{\text{int}}[0]\,s + c_{\text{ext}}^{(e)} > 0$.

\begin{note}
\textbf{Unbounded endpoints are allowed.} If a side of the interval is $\pm\infty$,
the corresponding boundary term in \eqref{eq:m1-closed} evaluates to $0$
\emph{provided} the integrand decays in that direction --- that is, provided
$\operatorname{sign}(\operatorname{Re}\alpha_s)$ matches the unbounded side. If it
does not, the integral genuinely diverges along the analytic path, and the
integrator returns \code{None}, which tells the caller to fall back to
\code{scipy.quad} (which handles $\pm\infty$ bounds natively via a coordinate
transform). Returning \code{None} is the universal ``I cannot do this one
analytically, try numerics'' signal throughout the engine.
\end{note}

This same function is also called directly by the \emph{grouped} Phase~J path
(Chapter~\ref{ch:grouped-phasej}) through its optional \code{pole\_tuples}
argument, which lets that path inject pre-summed residues. We mention it so the
shared interface is on your radar; the grouped path is the subject of the next
chapter.

\section{The $m=2$ polygon integrator}\label{sec:m2}

For $m=2$ the polytope is a convex \emph{polygon} in the plane $(s_0, s_1)$ ---
the intersection of the half-planes from the retardation constraints. The
integral \eqref{eq:exp-over-polytope} becomes $\iint_{\text{polygon}}
\exp(\alpha s_0 + \beta s_1)\,\dd A$. The engine evaluates it in four steps,
implemented across \code{\_integrate\_2d\_polygon\_modesum}
(\file{final\_integral.py:2175}) and its helpers.

\subsection{Step 1 --- build the polygon by clipping}

Start from a counter-clockwise bounding box of half-width \code{POLYGON\_BBOX\_CAP}
($=200$) and \emph{clip} it against each constraint half-plane in turn. The
clipping uses the \term{Sutherland--Hodgman} algorithm,
\code{\_clip\_polygon\_to\_halfplane} (\file{final\_integral.py:454}).

\begin{defn}
\textbf{Sutherland--Hodgman polygon clipping} is a classic computer-graphics
algorithm. It walks the edges of a polygon and, for each clipping half-plane,
keeps the vertices on the ``inside,'' inserts the intersection point wherever an
edge crosses the boundary, and drops the outside vertices. Applied once per
constraint, it produces the convex intersection of the box with all the
half-planes. Vertices lying exactly on a boundary are counted as inside --- a
measure-zero choice that does not affect the area integral.
\end{defn}

The driver \code{\_polygon\_from\_2d\_constraints} (\file{final\_integral.py:499})
substitutes the current external-time values into each constraint's affine offset
and clips:

\begin{lstlisting}
polygon = [(-bbox_cap, -bbox_cap), (bbox_cap, -bbox_cap),
           (bbox_cap, bbox_cap), (-bbox_cap, bbox_cap)]
for (a_int, a_ext, c0) in constraint_data:
    c_eff = float(c0) + sum(float(a_ext[j]) * float(free_ext_vals[j])
                            for j in range(len(a_ext)))
    polygon = _clip_polygon_to_halfplane(polygon, float(a_int[0]), float(a_int[1]), c_eff)
    if not polygon:
        return []
\end{lstlisting}

\noindent An empty result (the constraints are jointly infeasible at this
$\tau$) integrates to $0$.

\subsection{Steps 2--4 --- triangulate and map to the unit triangle}

Fan-triangulate the polygon from vertex $0$ (every triangle shares vertex $0$).
For each triangle, affine-map it to the \emph{unit triangle}
$\{0\le u, w;\ u+w\le 1\}$; the exponential integrand becomes
$\exp(\alpha_0 + p\,u + q\,w)$ for some complex $p, q$ read off the vertices. The
unit-triangle integral then has the closed form
\begin{equation}
  J(p, q) = \int_0^1\!\dd u\int_0^{1-u}\!\dd w\;\ee^{p u + q w}
   = \frac{1}{q}\left[\frac{\ee^{p} - \ee^{q}}{p - q} - \frac{\ee^{p} - 1}{p}\right],
  \label{eq:Jpq}
\end{equation}
implemented in \code{\_exp\_over\_unit\_triangle} (\file{final\_integral.py:356}).
The per-triangle wrapper \code{\_exp\_over\_triangle} (\file{final\_integral.py:406})
maps the vertices, calls $J$, and scales by the Jacobian
$|\det|\cdot\ee^{v_0\text{-term}}$.

\begin{gotcha}
\textbf{Equation \eqref{eq:Jpq} is a minefield of removable singularities.} The
denominators $q$, $p$, and $p-q$ each vanish in a different degenerate regime,
and the formula then takes the indeterminate $0/0$ form --- numerically, that is
catastrophic cancellation. The code guards \emph{all four} regimes with
4th-order Taylor expansions, triggered when any of $|p|, |q|, |p-q|$ drops below
\code{\_J\_TAYLOR\_EPS} $= 10^{-6}$. For instance the doubly-degenerate
$|p|,|q|\to 0$ case returns the moment expansion of the triangle directly:
\begin{lstlisting}
return (0.5
        + (p + q) / 6.0
        + (p*p)/24.0 + (q*q)/24.0 + (p*q)/24.0
        + (p**3 + q**3)/120.0
        + (p*p*q + p*q*q)/120.0)
\end{lstlisting}
the coefficients being the moments $\iint 1 = \tfrac12$, $\iint u = \iint w =
\tfrac16$, and so on.
\end{gotcha}

\subsection{The bilateral overflow guard}

\code{\_exp\_over\_triangle} returns \code{None} --- abort, fall back to scipy ---
whenever any of $|\operatorname{Re} p|$, $|\operatorname{Re} q|$,
$|\operatorname{Re}(v_0\text{-term})|$ exceeds \code{EXP\_REAL\_LIMIT} $= 600$.
Note this guard is \emph{bilateral}: it checks the magnitude on \emph{both} signs,
not just the overflow ($+$) side.

\begin{gotcha}
The bilateral guard is not an over-caution; it is a hard-won correctness fix. The
structural cancellation in \eqref{eq:Jpq}, $(\ee^p-\ee^q)/(p-q) -
(\ee^p-1)/p$, becomes dominated by floating-point noise when \emph{one} of
$\ee^p$, $\ee^q$ underflows to zero and the other does not. A 2026 optimisation
(commit 7f0bf05) relaxed the guard to one-sided, reasoning that the underflow
direction is harmless; in fact it introduced wrong-direction loop corrections in
the spike-reset $k=2$, $\ell=1$ case (which exercises the $m=2$ path heavily) and
was reverted. The lesson, recorded in the code comment at line~436: keep the
guard bilateral so the fragile-cancellation cases route to \code{scipy.nquad},
which handles the well-conditioned integral natively.
\end{gotcha}

\section{The $m\ge 3$ causal-poset integrator}\label{sec:m3}

For three or more integration variables the polygon picture no longer scales, and
the engine switches to a combinatorial decomposition built on the
\emph{structure} of the causal constraints. This is the most involved of the
three backends, \code{\_integrate\_nd\_polytope\_poset\_modesum}
(\file{final\_integral.py:1726}).

\subsection{The constraint set as a directed acyclic graph}

Read the constraint set $\{\Delta t_e > 0\}$ as a directed graph on the $m$
integration variables. Each inter-axis constraint $s_v - s_u > 0$ becomes a
directed edge $u\to v$ (read ``$u$ precedes $v$''); each single-variable
constraint becomes a scalar bound $s_v > c$ or $s_v < c$. Because all the
constraints come from causal orderings, this graph is acyclic --- a
\term{directed acyclic graph} (DAG), equivalently a \term{partial order} (poset).
This is extracted into the frozen \code{\_CausalPoset} dataclass
(\file{final\_integral.py:691}) by \code{\_extract\_causal\_poset}
(\file{final\_integral.py:720}), which classifies each constraint as an inter-axis
edge, a scalar lower bound, a scalar upper bound, or --- if it is none of these
clean forms (a mixed or shifted constraint) --- returns \code{None} to force the
scipy fallback.

\subsection{Linear extensions and the chain simplex}

The key geometric fact: the full polytope is the \emph{disjoint union of simplex
regions, one per linear extension of the poset}.

\begin{defn}
A \textbf{linear extension} of a partial order is a total order consistent with
it --- a way of laying all the variables in a single line $s_{\sigma(0)} \le
s_{\sigma(1)} \le \cdots$ that never violates a ``precedes'' relation. It is
exactly a topological sort of the DAG. Each linear extension corresponds to one
\textbf{chain simplex}
\[
  L \le s_{\sigma(0)} \le s_{\sigma(1)} \le \cdots \le s_{\sigma(m-1)} \le U_\sigma,
\]
and these simplices tile the polytope without overlap.
\end{defn}

The linear extensions are enumerated by \code{\_enumerate\_linear\_extensions}
(\file{final\_integral.py:806}), a recursive Kahn-style topological-sort generator
that yields every extension in deterministic order. The common scalar lower bound
$L$ is resolved by \code{\_causal\_poset\_consistent\_scalar\_lower}
(\file{final\_integral.py:857}); when no variable carries an explicit lower bound,
a physical-margin fallback sets $L = (\text{earliest external time}) -
\code{POSET\_PHYSICAL\_MARGIN}$ with \code{POSET\_PHYSICAL\_MARGIN} $= 50$.

\begin{gotcha}
Why the physical margin $50$ rather than the polygon's $200$? Because
\code{POLYGON\_BBOX\_CAP}$=200$ is too loose for the chain simplex: combined with
retarded poles ($|\operatorname{Re}\beta| > 3$ when three or four poles sum), the
factor $\ee^{\beta\cdot L}$ at $L=-200$ overflows the $600$ guard and bails to
scipy. The poset path uses the tighter margin $50$ --- several correlation times
beyond the earliest external time, where the retarded integrand is already
negligible --- so the analytic path actually completes. The comment block at
\file{final\_integral.py:290} records this trade-off in detail.
\end{gotcha}

\subsection{The chain-simplex closed form}

The nested exponential integral over one chain
$\{L\le s_1\le\cdots\le s_N\le U\}$ has a closed form obtained by integrating
inside-out. The reference implementation is \code{\_exp\_over\_chain\_simplex}
(\file{final\_integral.py:910}). Its derivation, quoted from the docstring:

\begin{quote}\small
Integrate inside-out. After each step the running integrand is a sum of terms
$C\cdot\exp(\beta\cdot s_{\text{outer}} + \text{const})$. Each inner integration
produces two new terms --- one ``upper-bound'' piece (whose $\beta$ merges with the
next-outer variable's coefficient) and one ``lower-bound'' piece (a numerical
prefactor $\ee^{\beta\cdot L}$ falls out). After $N$ steps the sum has $2^N$
constant terms; their sum is the integral.
\end{quote}

\noindent Two boundary behaviours are essential:

\begin{itemize}
  \item An \textbf{empty chain} ($U \le L$) returns $0$ --- the domain has zero
        measure. This is precisely the $\tau < 0$ situation where the chain-bottom
        lower bound (from a leaf at $t=0$) and the chain-top upper bound (from a
        leaf at $t=\tau<0$) cross. Causality makes the diagram vanish there, and
        the code returns exactly $0$.
  \item A \textbf{degenerate cumulative coefficient} ($|\beta| < \epsilon$) makes
        the closed form's $1/\beta$ singular; the function returns \code{None} so
        the caller switches to the polynomial-prefactor variant below. There is
        also the same bilateral $600$ overflow guard as in the polygon path.
\end{itemize}

Two variants extend this base case:

\begin{description}
  \item[\code{\_exp\_over\_chain\_simplex\_polynomial}
        (\file{final\_integral.py:1339})] When a level's cumulative $\beta$
        vanishes, the antiderivative is a \emph{polynomial} in $s$ rather than an
        exponential. This variant carries a polynomial (in the basis
        $(s-\text{lower})^k$) through the subsequent levels, using the closed form
        $\int (s-\text{lower})^k\,\ee^{\beta s}\,\dd s$. It returns \code{None}
        only on overflow, never on degenerate $\beta$ --- it is the engine's answer
        to the $1/\beta$ singularity.
  \item[\code{\_chain\_with\_intermediate\_uppers}
        (\file{final\_integral.py:1489})] The base chain assumes the only scalar
        upper bound sits at the top of the chain. When scalar uppers appear at
        \emph{intermediate} positions, this function computes a non-decreasing
        ``effective upper'' per position, groups them into levels, and enumerates
        monotone ``cut tuples'' marking where the chain crosses each lower upper
        bound, multiplying the chain-simplex value of each non-empty piece. It
        subsumes an older, cruder ``maximality bail.''
\end{description}

\subsection{The numba fast path}

For $m\ge 3$ the $2^N$ term expansion is the dominant cost, so it is also
reimplemented for speed. \code{\_exp\_over\_chain\_simplex\_numba\_core}
(\file{final\_integral.py:1219}) is a \code{numba}-compiled translation operating
on pre-allocated complex buffers.

\begin{defn}
\textbf{Numba} is a just-in-time (JIT) compiler for a subset of Python and NumPy.
You decorate a function with \code{@numba.njit} (``no-python JIT'') and the first
time it runs, numba compiles it to native machine code through LLVM --- typically
$10$--$100\times$ faster than the interpreted version for tight numeric loops. The
\code{cache=True} option persists the compiled code to disk so the compilation
cost is paid only once across runs.
\end{defn}

The numba core returns a \code{(status, value)} pair --- status $0$ for ok, $1$
for degenerate $\beta$, $2$ for overflow --- so the pure-Python caller can take the
same fallback branches it would for the reference implementation. The import is
guarded, and there is a master switch:

\begin{lstlisting}
try:
    import numba as _numba
    _HAVE_NUMBA = True
except ImportError:
    _HAVE_NUMBA = False
    _numba = None

USE_NUMBA_CHAIN_SIMPLEX = True
\end{lstlisting}

\noindent The dispatcher \code{\_exp\_over\_chain\_simplex\_fast}
(\file{final\_integral.py:1303}) prefers the numba core when it is available and
enabled, and otherwise uses the pure-Python reference --- which is guaranteed
bit-compatible on converged results. If numba is not installed, nothing breaks;
the engine is simply slower on $m\ge 3$ diagrams.

\section{The scipy fallback layer}\label{sec:scipy}

Whenever an analytic backend returns \code{None} --- degenerate, overflowing,
mixed-constraint, or unbounded-divergent --- the subset closure routes to a
numerical-quadrature fallback built on SciPy. The dispatcher
\code{\_integrate\_polytope} (\file{final\_integral.py:4372}) branches on $m$ and
forwards to \code{\_integrate\_1d\_polytope}, \code{\_integrate\_2d\_polytope}, or
\code{\_integrate\_nd\_polytope}, all built on \code{scipy.integrate.quad} (1D)
and \code{nquad} (nested multi-D). Real and imaginary parts are integrated
separately by \code{\_complex\_quad} (\file{final\_integral.py:4520}), since
\code{quad} works on real-valued integrands. The accuracy is tuned by the
module-level \code{QUAD\_OPTS = \{'limit': 200\}}.

This fallback layer has three subtle correctness requirements, each of which has
its own regression test:

\begin{description}
  \item[The Heaviside filter.] The polytope \emph{bounds} passed to \code{nquad}
        can overshoot the true causal region (the cap fallbacks and deferred
        cross-axis constraints are loose). For retarded poles the smooth
        propagator \emph{grows} for $\Delta t < 0$, so an overshooting box would
        let the integrand add a large spurious positive contribution.
        \code{\_make\_heaviside\_filtered\_integrand}
        (\file{final\_integral.py:4411}) wraps the integrand to return $0$
        whenever any $\Delta t_e \le 0$, enforcing the $\Theta(0)=0$ convention.
  \item[Deferred constraints.] In \code{\_integrate\_nd\_polytope}
        (\file{final\_integral.py:4805}) each axis's bound function must
        \emph{skip} any constraint that still couples to a more-inner axis ---
        otherwise the bound would depend on a variable not yet integrated. This is
        the deferred-constraint correctness handled by \code{\_make\_bound\_fn},
        guarded by the test
        \code{test\_phase\_J\_nd\_polytope\_preserves\_deferred\_constraints}.
  \item[The degenerate-empty sentinel.] \code{\_resolve\_1d\_bounds}
        (\file{final\_integral.py:4550}) returns a degenerate $(0.0, 0.0)$ --- not
        $(\infty, -\infty)$ --- when the constraints are infeasible. The reason is
        sharp: \code{scipy.quad(f, +inf, -inf)} returns the \emph{sign-flipped}
        full-line integral, which would silently poison any outer \code{nquad}.
        A zero-width interval correctly integrates to $0$.
\end{description}

\begin{gotcha}
A specific trap in \code{\_integrate\_2d\_polytope}: a constraint with both
$a_{\text{int}}[0]\approx 0$ and $a_{\text{int}}[1]\approx 0$ is a
\emph{pure-external} inequality (a condition on $\tau$ alone), not a bound on
$s_1$. Treating it as an $s_1$ bound sets a \code{pure\_s1\_found} flag that
skips the $\pm 200$ cap fallback and oversamples an unbounded axis --- the source
of roughly $12\%$ of the $k=4$ theory-versus-simulation overshoot. The code is
careful to distinguish the two.
\end{gotcha}

\section{Non-local kernels: noise sources and conductance vertices}\label{sec:kernels}

Two vertex types break the pure rational-propagator structure and need extra
integration variables. Both are recognised by importing the vertex classes
\code{NoiseSourceType} and \code{ConvVertexType} from \code{msrjd.core.vertices}.

\paragraph{\code{NoiseSourceType} (coloured-noise cumulants).} The source's legs
sit at \emph{independent} times coupled by a noise kernel $\kappa^{(n)}(\tau)$.
The engine adds one $\tau$ per noise vertex, routes one leg to
$\text{anchor} - \tau$, and substitutes the kernel $\kappa(\tau)$ into the
prefactor. Because $\kappa$ is generally \emph{not} a sum of exponentials --- a
Gaussian $\ee^{-\tau^2/2\sigma^2}$, say --- these diagrams cannot stay on the
analytic path and are forced through the symbolic-plus-\code{scipy.nquad} route.

\paragraph{\code{ConvVertexType} (conductance/synaptic kernels $g(\tau)$).} Same
scaffold, but with a rescue. If $g(\tau)$ \emph{single-exponential-extracts} ---
i.e. it is of the form $C\,\ee^{\lambda\tau}$ --- the engine can absorb it. The
extraction is \code{\_extract\_exp\_mode} (\file{final\_integral.py:203}), which
uses the log-derivative trick: $\lambda = \tfrac{\dd}{\dd\tau}\log g\big|_{\tau=0}$
and $C = g(0)$. A successfully extracted kernel becomes a synthetic
\textbf{pseudo-edge} $(C, \lambda)$ with constraint $\Delta t = +\tau > 0$, and
the diagram stays on the fast analytic path.

\begin{gotcha}
\code{\_extract\_exp\_mode} returns \code{None} for \emph{polynomial-prefactor}
kernels --- the classic example being the ``alpha'' synaptic kernel
$\tau/\tau_g^2\cdot\ee^{-\tau/\tau_g}$, for which $g(0)=0$. Such kernels need a
multi-mode decomposition the single-exponential extractor cannot provide, so they
fall back to the slow path. The single-exponential conductance kernel
$\tau_g\,\ee^{-\tau/\tau_g}$ is fine; the alpha kernel is not.
\end{gotcha}

\begin{gotcha}
The extra $\tau$ variables are capped to $\pm\code{TAU\_KERNEL\_CAP}$ ($=50$). This
is a quadrature band-aid for the slow path: without the cap, \code{scipy.quad}
over $(\text{retard}_L, +\infty)$ lets the standard $\tan$-substitution compress a
Gaussian kernel's peak against the boundary, where the adaptive sampler
intermittently misses it and produces spurious spikes. Conductance $\tau$ uses a
lower bound of $0$ (causality); noise $\tau$ uses the symmetric $\pm$ cap.
\end{gotcha}

\section{Batching and parallelism}\label{sec:parallel}

The orchestration layer (\file{pipeline.py}) builds, on top of the per-diagram
callables, two batched evaluators for sweeping a $\tau$ grid:

\begin{itemize}
  \item \code{total\_C\_batch(tau\_points, parallel=True, ...)}
        (\file{pipeline.py:374}, a closure inside \code{compute\_correction\_td})
        --- evaluates the full summed correlator over a grid.
  \item \code{eval\_per\_diagram\_batch(...)} (\file{pipeline.py:502}) --- returns
        the full \code{(n\_tau, n\_diagram)} grid of per-diagram contributions, for
        diagnostics.
\end{itemize}

Both can fan out over a pool of worker processes. The parallelism uses
\code{multiprocessing.Pool} with the \code{fork} start method --- and that choice
is forced, for a reason that also makes it dangerous.

\begin{defn}
\textbf{Fork vs.\ spawn.} These are the two ways the \code{multiprocessing} module
creates worker processes. \code{fork} (POSIX only) copies the parent's entire
memory image into each child --- the children inherit every live object for free.
\code{spawn} starts a fresh Python interpreter and \emph{pickles} the arguments
across the process boundary.
\end{defn}

Fork is \emph{mandatory} here because the per-diagram \code{contribution}
callables are nested closures (defined inside
\code{integrate\_diagram.<locals>.contribution}), and the standard-library
\code{pickle} cannot serialise a nested closure. With \code{fork}, the workers
inherit those closures through the copied memory image; a module-level
\code{\_WORKER\_STATE} dict carries them in. The module-level worker entry points
\code{\_worker\_eval\_total\_C} (\file{pipeline.py:174}) and
\code{\_worker\_eval\_one\_diagram} (\file{pipeline.py:186}) are themselves
picklable (they are top-level functions) and simply look up the inherited
closures.

There are two parallel strategies, auto-chosen by batch shape: per-$\tau$
parallelism (each worker evaluates the full \code{total\_C} at one $\tau$) when
there are at least \code{n\_workers} grid points, and per-$(\tau, \text{diagram})$
nested parallelism otherwise --- the latter aggregates results in ascending
diagram order so that floating-point accumulation is bit-identical to the serial
path.

\begin{gotcha}
\textbf{Fork can crash the OS inside a notebook.} Forking a process \emph{after}
the Cocoa/BLAS libraries have initialised --- which is the state of any macOS
Jupyter kernel --- has hard-crashed the development machine twice. The environment
variable \code{OBJC\_DISABLE\_INITIALIZE\_FORK\_SAFETY=YES} makes it \emph{worse}:
it removes the guard rail, not the hazard. Phase~J therefore consults
\code{\_fork\_unsafe\_in\_notebook} (\file{pipeline.py:146}) --- which is true iff
the start method is \code{fork} \emph{and} the platform is macOS \emph{and} the
process is running inside a ZMQ/Jupyter kernel --- and silently degrades to serial
evaluation with a one-time warning. The bit-identity tests
\code{test\_phase\_J\_total\_C\_batch\_*} pass under \emph{pytest}, which is a plain
interpreter and never a ZMQ kernel, so they certify numerical determinism but
\emph{not} notebook crash-safety. The lesson, learned twice: in a notebook on a
Mac, keep \code{parallel=False}.
\end{gotcha}

\begin{gotcha}
Windows has no \code{fork}, and the per-diagram closures are unpicklable, so
\code{total\_C\_batch(parallel=True)} raises \code{ValueError: cannot find context
for 'fork'} on Windows. The parallel batch path is POSIX-only; on Windows, run
serially.
\end{gotcha}

\section{$\delta$-spike discretisation}\label{sec:deltaspike}

Recall from \S\ref{sec:polytope} that a $\delta$-edge equation with no
integration variable to solve for becomes a constraint among external times --- a
distributional shot-noise spike $\delta(a\cdot\tau + c)$. These are collected
into the result's \code{delta\_contributions} list rather than folded into the
smooth callable, because a $\delta$-function has no value to evaluate pointwise;
it must be deposited onto a grid.

\code{eval\_delta\_contributions\_on\_tau\_grid} (\file{final\_integral.py:3828})
does this for a 1D $\tau$ grid. For each contribution with equality
$a\cdot x + c = 0$ it solves for the firing lag $\tau_{\text{fire}} =
-c'/a_{\text{vary}}$, evaluates the (compiled) coefficient callable there, checks
the retardation half-spaces, and deposits $\text{coeff}/|a_{\text{vary}}|/\dd\tau$
into the nearest bin --- the $1/\dd\tau$ being the discrete stand-in for the
$\delta$'s unit area.

\begin{gotcha}
There is a degenerate case the code handles explicitly: when $a_{\text{vary}}
\approx 0$ \emph{and} $c'\approx 0$, the $\delta$ fires along the \emph{whole}
slice rather than at a single point. The code then deposits the
\emph{continuous} coefficient $\text{coeff}_{\text{fc}}(\tau)$ with \emph{no}
$1/\dd\tau$ divisor --- treating it as a smooth contribution along the slice, not a
point spike.
\end{gotcha}

The 2D analogue \code{eval\_delta\_contributions\_on\_2d\_grid}
(\file{final\_integral.py:4004}) sweeps the $\delta$-line
$a_1\tau_1 + a_2\tau_2 + c = 0$ along whichever axis carries the larger
coefficient, depositing into the nearest cell.

\section{External tools used}\label{sec:tools}

This subsystem straddles SageMath and the scientific-Python stack. Each library
appears for a specific, narrow reason; here is the inventory.

\subsection{SageMath (\code{sage.all})}

SageMath is a large open-source mathematics system that wraps dozens of
specialised libraries (Maxima for symbolic algebra, PARI, GMP, and more) behind a
unified Python API. Phase~J uses it only as the \emph{authoring and extraction}
layer --- the diagram's propagator entries, the $\delta$-coefficients, and the
prefactor all arrive as Sage symbolic expressions. The relevant objects:

\begin{description}
  \item[\code{SR} --- the Symbolic Ring.] \code{SR.var('x')} makes a symbolic
        variable; \code{SR(expr)} coerces a value into a symbolic expression.
        Symbolic expressions support \code{.subs(dict)}, \code{.expand()},
        \code{.diff()}, \code{.coefficient(v)}, and \code{.variables()}.
  \item[\code{CDF} --- the Complex Double Field.] Machine-precision complex
        numbers. \code{complex(CDF(SR(expr)))} is the standard idiom for turning a
        fully-substituted symbolic scalar into a Python \code{complex}.
  \item[\code{fast\_callable(expr, vars, domain=CDF)}.] Sage's JIT compiler
        \emph{for symbolic expressions}. It walks the expression tree once and
        produces a fast callable (backed by a bytecode interpreter) that evaluates
        \code{expr} numerically. Phase~J uses it to compile the residual symbolic
        integrand for the \code{scipy.nquad} fallback path only.
  \item[\code{solve} (imported as \code{sage\_solve}).] The symbolic equation
        solver, used on the $\delta$-edge equations $\Delta t_e == 0$ to solve for
        an integration variable.
\end{description}

\noindent The exact import (\file{final\_integral.py:83}):

\begin{lstlisting}
from sage.all import SR, fast_callable, CDF, solve as sage_solve
\end{lstlisting}

The discipline is: extract the numerical pole/residue data \emph{once}, solve the
$\delta$ equations, read linear coefficients off the constraint expressions --- and
then never touch Sage again. The hot numerical loops are pure Python
\code{complex} and \code{cmath}, because Sage is too slow per call.

\subsection{\code{cmath} / \code{math} (standard library)}

\code{cmath.exp} is the complex exponential at the core of every analytic
integrator's inner loop; \code{math.inf}, \code{math.isinf}, and
\code{math.factorial} appear in the bound resolution and the polynomial-prefactor
chain. These are imported \emph{locally} inside the hot functions (e.g.
\code{import cmath}) so the module-level namespace stays clean.

\subsection{NumPy and Numba}

\term{NumPy} (\code{import numpy as np}, \file{final\_integral.py:1041}) plays two
roles: it backs the numba chain-simplex kernel, which works on pre-allocated
\code{np.complex128} buffers, and it provides the output arrays (and
\code{np.argmin}/\code{np.abs}) for the $\delta$-spike grid discretisers.
\term{Numba} is the JIT compiler for the chain-simplex inner loop covered in
\S\ref{sec:m3}; its import is guarded so the engine runs without it.

\subsection{mpmath}

\term{mpmath} is a pure-Python arbitrary-precision floating-point library;
\code{mpmath.mpc} is a complex number with a configurable number of decimal
digits (\code{mp.dps}). Phase~J uses it in two \emph{experimental,
off-by-default} precision-rescue paths for the $m\ge 3$ chain simplex, where the
$2^N$ cancellation can lose precision on close-paired poles. Both
\code{USE\_CHAIN\_SIMPLEX\_PRECISION\_FIX} and
\code{USE\_POSET\_MPMATH\_ACCUMULATION} default to \code{False} --- they did not
actually cure the targeted bug (see \S\ref{sec:knownbugs}).

\subsection{SciPy}

\term{SciPy}'s \code{integrate} submodule provides adaptive numerical quadrature:
\code{quad(f, a, b)} for adaptive 1D integration (Gauss--Kronrod / QUADPACK,
handling $\pm\infty$ bounds natively), and \code{nquad(f, [bounds...])} for nested
multi-dimensional quadrature where each axis's bounds may be a callable of the
outer variables. It is the engine's universal fallback, reached whenever an
analytic integrator returns \code{None}.

\subsection{What is \emph{not} used}

NetworkX, nauty, and SymPy do \emph{not} appear in this subsystem. Graph
structure is handled through SageMath's own graph objects
(\code{D.vertices()}, \code{D.edges()}, \code{D.num\_edges()}); nauty (graph
canonicalisation) lives in the upstream enumeration subsystem; and all symbolic
work here is Sage's \code{SR}, never SymPy.

\section{Diagnostics and tuning knobs}\label{sec:diagnostics}

The engine exposes a handful of module-level knobs and a runtime-counter
mechanism, all so you can see \emph{which} path actually fired and tune the
trade-offs.

\begin{description}
  \item[\code{\_RUNTIME\_COUNTERS} (\file{final\_integral.py:315}).] A diagnostic
        dict counting, per $m$, how often each analytic path was \emph{attempted},
        how often it \emph{returned \code{None}}, and how often the run fell
        through to \code{scipy.nquad}. Call \code{\_reset\_runtime\_counters()}
        before a timed run, then read the dict after. These are \emph{runtime}
        observations; the \emph{intent} of each subset (which path it was
        \emph{set up} to take) is recorded separately in
        \code{subset\_diagnostics['evaluator']}.
  \item[\code{QUAD\_OPTS = \{'limit': 200\}}.] The accuracy knob for the scipy
        fallback's quadrature.
  \item[\code{TAU\_KERNEL\_CAP = 50.0}.] The cap on the extra non-local-kernel
        $\tau$ variables (\S\ref{sec:kernels}).
  \item[\code{POLYGON\_BBOX\_CAP = 200.0} vs.\ \code{POSET\_PHYSICAL\_MARGIN = 50.0}.]
        The bounding-box half-widths for the $m=2$ and $m\ge 3$ paths respectively
        (\S\ref{sec:m2}, \S\ref{sec:m3}).
  \item[The \code{USE\_*} flags.] \code{USE\_NUMBA\_CHAIN\_SIMPLEX} (default
        \code{True}) toggles the numba fast path; \code{USE\_POLYGON\_M2\_INTEGRATOR}
        and \code{USE\_1D\_INTEGRATOR} toggle the analytic $m=2$ and $m=1$ paths;
        and the precision-rescue flags
        \code{USE\_CHAIN\_SIMPLEX\_PRECISION\_FIX},
        \code{USE\_POSET\_MPMATH\_ACCUMULATION}, and
        \code{USE\_POSET\_CAP\_MATCH\_SCIPY} all default \code{False}.
\end{description}

\section{Known precision caveats and open bugs}\label{sec:knownbugs}

Honesty about the limits of the engine matters more than a clean narrative. There
is one genuinely open bug and a cluster of near-misses around it.

\begin{gotcha}
\textbf{The $m\ge 3$ chain-simplex close-pole overestimate (OPEN).} On
close-paired poles --- e.g.\ the spike-reset poles $0.329\ii$ and $0.351\ii$, a
difference of only $0.022\ii$ --- the $2^N$ closed form suffers
cancellation-driven precision loss. Empirically this shows up as roughly a
$4\times$ aggregate overestimate of the one-loop value at spike-reset $k=2$,
$\ell=1$. Two attempted fixes are \emph{both off by default} because neither
actually cured it: \code{USE\_CHAIN\_SIMPLEX\_PRECISION\_FIX} (an mpmath dispatch
on close poles) and \code{USE\_POSET\_MPMATH\_ACCUMULATION} (mpmath outer
accumulation, which improved the cancellation factor only $\sim 21\times$). The
audit (\file{docs/m\_ge3\_precision\_bug\_audit.md}) now points at a
\emph{bookkeeping difference in pole-tuple construction between the grouped and
per-diagram paths} --- by linearity they should agree, so the $4\times$ discrepancy
points to a construction bug, not an arithmetic-precision one. This is
unresolved.
\end{gotcha}

The related near-miss flag \code{USE\_POSET\_CAP\_MATCH\_SCIPY} was meant to align
the analytic poset's unbounded-below fallback ($L = \text{earliest\_ext} - 50$)
with the scipy \code{OUTER\_CAP} of $200$, for marginally-stable theories where
the integral is genuinely cap-dependent. The code comment records that aligning
the caps does \emph{not} bring the per-diagram path into agreement with the
grouped path --- confirming the dominant error lies elsewhere. It too defaults
\code{False}.

\section{The loop-subgraph stub}\label{sec:subgraph}

For completeness: \file{subgraph.py} is a near-stub for a route that was designed
but never built. Its \code{identify\_loop\_subgraphs} (\file{subgraph.py:66})
returns \code{[]} for tree-level diagrams (empty \code{free\_freqs}) and
\emph{raises \code{NotImplementedError}} otherwise. The frequency-domain
loop-reduction algorithm this file was meant to enable was superseded: the
current \code{integrate\_diagram} handles loops directly in the time domain,
treating every internal vertex as an integration variable regardless of loop
order. The module is never exercised on the live path; its docstring documents
only the \emph{planned} algorithm. If you go looking for ``where the loop
reduction happens,'' the answer is that there is no separate loop reduction ---
the DAG-and-polytope machinery of this chapter \emph{is} the loop handling.

\begin{note}
The same point bears on a stale docstring you may notice while reading the code:
\file{pipeline.py}'s module header and the \code{compute\_correction\_td}
docstring both say loop diagrams are \emph{skipped}. They are not. The live code
evaluates every loop order through \code{integrate\_diagram}; the inline comment at
\file{pipeline.py:291} explains why (the enumerator produces DAGs, so tree and
loop share the same vertex-time integration algorithm). Trust the code, not the
header.
\end{note}

\section{What you just learned}

Phase~J turns each typed diagram into a number by recognising that the retarded
propagator is a finite sum of decaying exponentials \eqref{eq:GR-decomp}, so the
diagram integral \eqref{eq:diagram-integral} is a sum --- over $\delta$-subsets
and pole tuples --- of pure exponentials integrated over a causal polytope
\eqref{eq:exp-over-polytope}. Three analytic backends evaluate that polytope
integral in closed form, dispatched by the surviving dimension $m$: an interval
formula for $m=1$ (\S\ref{sec:m1}), Sutherland--Hodgman clipping plus a
unit-triangle formula for $m=2$ (\S\ref{sec:m2}), and a causal-poset
linear-extension decomposition into chain simplices for $m\ge 3$
(\S\ref{sec:m3}). A \code{scipy} quadrature layer catches every degeneracy
(\S\ref{sec:scipy}), non-local noise and conductance kernels add extra
$\tau$ variables (\S\ref{sec:kernels}), distributional spikes are deposited onto
the grid separately (\S\ref{sec:deltaspike}), and a fork-based batch layer ---
with a hard-won notebook safety guard --- sweeps the $\tau$ grid
(\S\ref{sec:parallel}).

The natural next step is Chapter~\ref{ch:grouped-phasej}, the \emph{grouped}
Phase~J path. It reuses the very same analytic primitives ---
\code{\_integrate\_1d\_polytope\_modesum}, \code{\_integrate\_2d\_polygon\_modesum},
and the poset integrator --- but sums residues \emph{across} diagrams before
integrating, trading per-diagram independence for a reduction in the pole-tuple
count. The discrepancy between that path and this one is, as
\S\ref{sec:knownbugs} noted, the current prime suspect for the $m\ge 3$ open bug.
